Rozložení mezd podle percentilů (P25, P50, P75) a~mezikvartilového rozpětí (\ac{IQR}) ve vybraných odvětvích ČR (\ac{ISPV}) odhaluje strukturální nerovnost uvnitř odvětví: v~maloobchodě (G) a~pohostinství (I) je \ac{IQR} úzké a~P25 těsně nad minimální mzdou, zatímco v~ICT (J) a~profesních službách je \ac{IQR} výrazně širší s~P75 vysoko nad mediánem.
Úzké \ac{IQR} u~nízkopříjmových odvětví svědčí o~komprimaci způsobené zákonným minimem, nikoli kolektivně sjednanými tarifními úrovněmi: bez \ac{KS} je minimální mzda jedinou pracovněprávní zárukou a~trh nevytváří tlak na její překonání.
Výrazné \ac{IQR} v~sektoru ICT dokládá, že v~odvětvích s~individuálním vyjednáváním dochází k~mzdové polarizaci uvnitř profesní skupiny --- výsledek, který \ac{KS} s~mzdovými tabulkami zabraňují tím, že nastavují transparentní tarifní úrovně pro jednotlivé kvalifikační stupně.
Porovnání odvětvových percentilních profilů tak potvrzuje, že chybějící sektorová \ac{KS} v~CZ jsou přímo spojena s~vyšší vnitroanalytickou mzdovou nerovností v~odvětvích bez kolektivního pokrytí.

Rozložení mezd podle percentilů (P25, P50, P75) a~mezikvartilového rozpětí (\ac{IQR}) ve vybraných odvětvích ČR (\ac{ISPV}) odhaluje strukturální nerovnost uvnitř odvětví: v~maloobchodě (G) a~pohostinství (I) je \ac{IQR} úzké a~P25 těsně nad minimální mzdou, zatímco v~ICT (J) a~profesních službách je \ac{IQR} výrazně širší s~P75 vysoko nad mediánem.
Úzké \ac{IQR} u~nízkopříjmových odvětví svědčí o~komprimaci způsobené zákonným minimem, nikoli kolektivně sjednanými tarifními úrovněmi: bez \ac{KS} je minimální mzda jedinou pracovněprávní zárukou a~trh nevytváří tlak na její překonání.
Výrazné \ac{IQR} v~sektoru ICT dokládá, že v~odvětvích s~individuálním vyjednáváním dochází k~mzdové polarizaci uvnitř profesní skupiny --- výsledek, který \ac{KS} s~mzdovými tabulkami zabraňují tím, že nastavují transparentní tarifní úrovně pro jednotlivé kvalifikační stupně.
Porovnání odvětvových percentilních profilů tak potvrzuje, že chybějící sektorová \ac{KS} v~CZ jsou přímo spojena s~vyšší vnitroanalytickou mzdovou nerovností v~odvětvích bez kolektivního pokrytí.

Rozložení mezd podle percentilů (P25, P50, P75) a~mezikvartilového rozpětí (\ac{IQR}) ve vybraných odvětvích ČR (\ac{ISPV}) odhaluje strukturální nerovnost uvnitř odvětví: v~maloobchodě (G) a~pohostinství (I) je \ac{IQR} úzké a~P25 těsně nad minimální mzdou, zatímco v~ICT (J) a~profesních službách je \ac{IQR} výrazně širší s~P75 vysoko nad mediánem.
Úzké \ac{IQR} u~nízkopříjmových odvětví svědčí o~komprimaci způsobené zákonným minimem, nikoli kolektivně sjednanými tarifními úrovněmi: bez \ac{KS} je minimální mzda jedinou pracovněprávní zárukou a~trh nevytváří tlak na její překonání.
Výrazné \ac{IQR} v~sektoru ICT dokládá, že v~odvětvích s~individuálním vyjednáváním dochází k~mzdové polarizaci uvnitř profesní skupiny --- výsledek, který \ac{KS} s~mzdovými tabulkami zabraňují tím, že nastavují transparentní tarifní úrovně pro jednotlivé kvalifikační stupně.
Porovnání odvětvových percentilních profilů tak potvrzuje, že chybějící sektorová \ac{KS} v~CZ jsou přímo spojena s~vyšší vnitroanalytickou mzdovou nerovností v~odvětvích bez kolektivního pokrytí.

Rozložení mezd podle percentilů (P25, P50, P75) a~mezikvartilového rozpětí (\ac{IQR}) ve vybraných odvětvích ČR (\ac{ISPV}) odhaluje strukturální nerovnost uvnitř odvětví: v~maloobchodě (G) a~pohostinství (I) je \ac{IQR} úzké a~P25 těsně nad minimální mzdou, zatímco v~ICT (J) a~profesních službách je \ac{IQR} výrazně širší s~P75 vysoko nad mediánem.
Úzké \ac{IQR} u~nízkopříjmových odvětví svědčí o~komprimaci způsobené zákonným minimem, nikoli kolektivně sjednanými tarifními úrovněmi: bez \ac{KS} je minimální mzda jedinou pracovněprávní zárukou a~trh nevytváří tlak na její překonání.
Výrazné \ac{IQR} v~sektoru ICT dokládá, že v~odvětvích s~individuálním vyjednáváním dochází k~mzdové polarizaci uvnitř profesní skupiny --- výsledek, který \ac{KS} s~mzdovými tabulkami zabraňují tím, že nastavují transparentní tarifní úrovně pro jednotlivé kvalifikační stupně.
Porovnání odvětvových percentilních profilů tak potvrzuje, že chybějící sektorová \ac{KS} v~CZ jsou přímo spojena s~vyšší vnitroanalytickou mzdovou nerovností v~odvětvích bez kolektivního pokrytí.

\input{texparts/python/rscp_sector_percentiles}



